\Needspace{15\baselineskip}
\fancyhead[L]{\small GAUSS TO VON ZACH.}
\begin{center}
GAUSS TO VON ZACH.
\end{center}

\begin{center}
Monatliche Correspondenz. Vol. XIV. Pp. 377–382. October 1806.
\end{center}

\begin{flushright}
Brunswick, 25 August 1806.
\end{flushright}

\par\indent --- --- --- I enclose the ephemeris for Juno. It has been ready for some time, but I kept hoping I could also include the one for \emph{Ceres}. I still have no entirely reliable observations of Ceres from around the time of
\SourceLeaf{285}%
opposition. --- If I do not receive good observations of it soon, I will have to improve the elements for the time being with \textsc{Pasquich}’s observations, which in his own judgment are not very reliable. \textsc{Harding}’s observations and a couple that \textsc{Bode} sent me were made somewhat farther from the time of opposition. Might it not also have been observed in Milan?

\par\indent These are the elements from which the ephemeris for \emph{Juno} was calculated:

\begin{center}
VI. \emph{Elements of Juno.}
\end{center}

\Needspace{18\baselineskip}
\begin{center}
Epoch of longitude on the Seeberg meridian
\end{center}
\[
\left\{\begin{array}{ll}
1805. & 42^\circ\,35'\,7''3\\
1806. & 125\;11\;20.1\\
1807. & 207\;47\;33.0
\end{array}\right.
\]
\begin{center}\begingroup
\fontsize{9}{11}\selectfont
\setlength{\tabcolsep}{3pt}\renewcommand{\arraystretch}{1.12}
\begin{tabular}{ll}
daily mean tropical motion  & $814''7201$ \\
Longitude of aphelion in 1806 .\ .\ . & $233^\circ\,17'\,1''$ \\
Longitude of the ascending node in 1806  & $171^\circ\,4'\,57''4$ \\
Eccentricity .\ .\ .\ .\ .\ .\ .\ .\ . & $0.254944$ \\
Logarithm of the semiaxis .\ .\ .\ . & $0.4260480$ \\
Inclination of the orbit .\ .\ .\ .\ .\ .\ . & $13^\circ\,3'\,28''4$ \\
\end{tabular}\endgroup\end{center}

\par\indent Prof. \textsc{Harding} writes to me that he has already finished most of both charts, and that feeling unwell is holding up only the comparison with the sky itself. He hopes to be able to send them to you soon.

\Needspace{15\baselineskip}
\par\indent I did not include the brightness of \emph{Juno} next year in my ephemeris because this time I calculated it by a new and particularly convenient method, in which the planet’s distance from the Sun does not enter into the calculation at all. For an overview, I will nevertheless give three calculated brightness values, which show how extraordinarily difficult it will be to observe Juno next year:

\begin{center}\begingroup
\fontsize{9}{11}\selectfont
\setlength{\tabcolsep}{3pt}\renewcommand{\arraystretch}{1.12}
\begin{tabular}{ll}
1807 January 12 .\ . & $0.00661,$ \\
May 18 .\ . & $0.01600,$ \\
Sept. 21 .\ . & $0.00709$ \\
\end{tabular}\endgroup\end{center}

\par\indent This year in March the greatest brightness was still $0.03983$ (M.\,C., vol.~XI, \href{https://www.digitale-sammlungen.de/en/view/bsb10538603?page=492}{p.~478} [\hyperlink{native-270}{p.~265} of this volume]), so even at its brightest in 1807, Juno will have less than half that brightness. --- The new method just mentioned shortens the calculation so much that I can fit 15 complete calculations on an ordinary octavo page if the positions of the Sun are taken as already given.

\par\indent [The geocentric course of Juno from 12 January to 21 September 1807 is shown at three-day intervals for 12 o’clock, midnight at Seeberg, with right ascension and declination given to the minute and the distance from the Earth to three decimal places.]

